\documentclass{exam}
\usepackage{../../commonheader}

\discnumber{9}
\title{Inheritance, Iterators, \& Generators}
\date{October 19, 2026 -- October 23, 2026}

% Exercises and lessons are imported from the topic bank when possible.
\begin{document}
\maketitle

\begin{meta}
\small
\textbf{Example Timeline}
\begin{itemize}
    \item \textbf{This is a menu of problems, not a requirement to finish the worksheet.} Individual suggested times do not add up to one hour. Choose or skip problems based on your students, and give each mini-lecture immediately before its section.
    \item Check-in and icebreaker [5 min]. Ask students for an example of an ``is-a'' relationship (for instance, a checking account is an account).
    \item Inheritance: Q1 (WWPD with \lstinline{Elphaba} and \lstinline{Glinda}), then Q2 (\lstinline{ForceWielder}) or Q3 (\lstinline{TeamBaller}) [20 min]. Q1 builds the lookup chain and \lstinline{type(self)}; Q2 adds \lstinline{super()} and class attributes; Q3 practices overriding and calling the base class's method.
    \item Iterators: Q4 (iterator versus iterable), then Q5 and Q6 (WWPD) [10 min]. Q7 is a challenge: students work backward from the printed output by tracking iterator positions.
    \item Generators: Q8 (WWPD) [8 min], then Q9 (\lstinline{accumulate}) for generator syntax with an infinite iterator [7 min].
    \item Challenge: Q10 (\lstinline{fruitOptions}) is a recursive generator. Use it if students are comfortable with \lstinline{yield} in a \lstinline{for} loop over a recursive call.
    \item Wrap-up and questions [5 min]. Ask students to explain the difference between calling a function and calling a generator function.
\end{itemize}
\textbf{Teaching Goals}
\begin{itemize}
    \item A subclass inherits from a base class without copying its attributes. Attribute lookup checks the instance, then the class, then each base class up the chain, and only ever goes up.
    \item Overriding replaces an inherited method in the subclass; \lstinline{super()} (or \lstinline{Base.method(self, ...)}) lets the new method call the old one, and \lstinline{self} remains the subclass instance.
    \item Lookup and assignment differ: assignment through an object writes to that object, so \lstinline{self.x = ...} makes an instance attribute and \lstinline{Class.x = ...} changes the class attribute.
    \item \lstinline{iter} creates an iterator from an iterable, and \lstinline{next} advances it. Two names for one iterator share a position; two calls to \lstinline{iter} on a list create two iterators, but \lstinline{iter} on an iterator returns that same iterator.
    \item \lstinline{map}, \lstinline{filter}, \lstinline{zip}, and generators are lazy: nothing is computed until a consumer asks for the next element. \lstinline{list} comprehensions are eager.
    \item A generator function returns a generator without running its body. The body runs on \lstinline{next}, pauses at \lstinline{yield}, and resumes where it left off; the body finishing raises \lstinline{StopIteration}.
    \item In recursive generators, use \lstinline{yield from} (or a \lstinline{for} loop) to pass along the values from a recursive call, and \lstinline{yield} to produce your own values.
\end{itemize}
\end{meta}

\section{Inheritance}
\begingroup
\subimport{../../topics/oop/text/}{inheritance.tex}
\endgroup
\begin{questions}
    \subimport{../../topics/oop/easy/wwpd/}{wicked.tex}

    \subimport{../../topics/oop/easy/}{star-wars.tex}

    \subimport{../../topics/oop/easy/}{team-baller-fa21.tex}
\end{questions}

\section{Iterators}
\begingroup
\subimport{../../topics/iterators/text/}{iterators.tex}
\endgroup
\begin{questions}
    \setcounter{question}{3}
    \subimport{../../topics/iterators/easy/}{iterator-vs-iterable.tex}

    \subimport{../../topics/iterators/easy/}{wwpd.tex}

    \subimport{../../topics/generators/ExamRev-Sp25/}{numbers_and_letters.tex}

    \subimport{../../topics/generators/ExamRev-Sp25/}{year.tex}
\end{questions}

\section{Generators}
\begingroup
\subimport{../../topics/generators/text/}{generators.tex}
\endgroup
\begin{questions}
    \setcounter{question}{7}
    \subimport{../../topics/generators/easy/}{foo.tex}

    \subimport{../../topics/generators/easy/}{accumulate.tex}

    \subimport{../../topics/generators/ExamRev-Sp25/}{fruit_options.tex}
\end{questions}

\end{document}
